\documentclass{article}

% Configuración de idioma (español para la mayor parte del texto)
\usepackage[spanish, english]{babel}
\usepackage[utf8]{inputenc}

% Márgenes cómodos para leer
\usepackage[letterpaper, top=2cm, bottom=2cm, left=2.5cm, right=2.5cm]{geometry}

% Para usar tildes y caracteres directamente
\usepackage[T1]{fontenc}

% Paquetes útiles
\usepackage{amsmath}
\usepackage{graphicx}
\usepackage[colorlinks=true, allcolors=blue]{hyperref}

\title{Guion de presentación: KnowGraph}
\author{Defensa de proyecto}
\date{} % sin fecha

\begin{document}
\maketitle

\section*{Introducción}

\subsection*{Diapositiva: Título}
\textit{(Sin hablar, se muestra la portada mientras el tribunal se acomoda)}

\subsection*{Diapositiva: Contenidos}
\textit{(Iniciar tras pasar a esta diapositiva)}

Buenos días/tardes. Mi nombre es Valentino Russo y hoy presento mi proyecto final de ingeniería, el titulo de la aplicación es \textbf{KnowGraph}.

En esta exposición recorreré primero la motivación y el contexto, los conceptos necesarios para entender la herramienta, una demostración de su funcionamiento, los detalles técnicos para los desarrolladores y, finalmente, una reflexión sobre sus puntos fuertes y débiles, junto con posibles mejoras. La idea es que tanto quien viene del mundo técnico como quien no, pueda seguir el hilo sin dificultad.

\section*{Sección 1: Motivación y proyectos similares}

\subsection*{Diapositiva: Motivación}
El punto de partida es una necesidad real: explorar grandes volúmenes de datos interconectados sin ser un experto en grafos de conocimiento ni en SPARQL. Queríamos una herramienta que permitiera consultar, visualizar y transformar datos con una curva de aprendizaje baja, pero sin sacrificar potencia. KnowGraph integra ejecución de consultas, visualización interactiva (tanto en grafo como mediante reducción dimensional) y carga de datos tabulares, todo en un solo entorno.

\subsection*{Diapositiva: Proyectos y herramientas relevantes}
No partimos de cero. Existen referentes como Protégé, centrado en ontologías; Stardog, una base de datos semántica empresarial; y servicios como Wikidata Query Creator o DBpedia, que ofrecen consultas visuales sobre grandes grafos públicos. Sin embargo, estas herramientas suelen estar orientadas a un perfil muy técnico o ligadas a un único proveedor. KnowGraph busca ser más accesible, combinando las ventajas de un entorno local con la capacidad de conectarse a servicios externos.

\section*{Sección 2: Conocimientos necesarios}

\subsection*{Diapositiva: Conocimientos necesarios para utilizar la aplicación}
Aunque la aplicación intenta ocultar la complejidad, hay cuatro pilares conceptuales que conviene tener presentes: qué es un grafo, cómo se modela conocimiento con RDF, cómo se interroga ese conocimiento con SPARQL y cómo Docker facilita el despliegue. Los revisaremos brevemente.

\subsection*{Diapositiva: Grafos y representación de conocimiento}
Un grafo es simplemente un conjunto de nodos (entidades) y aristas (relaciones). Dependiendo de si la relación tiene dirección, el grafo puede ser dirigido o no; además, se pueden añadir etiquetas que describen el significado de cada relación. Esta estructura es muy natural para modelar cualquier dominio: personas, ciudades, genes, transacciones…

\subsection*{Diapositiva: Grafos de conocimiento}
Cuando a ese grafo le añadimos un modelo semántico explícito (es decir, definimos qué significan realmente las relaciones y las entidades), hablamos de grafos de conocimiento. Son la columna vertebral de la Web Semántica y del Linked Data, y permiten integrar información de fuentes heterogéneas razonando sobre ella.

\subsection*{Diapositiva: RDF y su conexión con los grafos}
El estándar técnico que materializa esto es RDF, Resource Description Framework. Define un modelo basado en triplas: sujeto, predicado, objeto. Cada tripla es una arista dirigida y etiquetada. RDF es la base sobre la que luego se construyen las consultas. Es interoperable y procesable automáticamente.

\subsection*{Diapositiva: SPARQL}
Para extraer información de esos grafos RDF utilizamos SPARQL, el lenguaje de consulta estándar del W3C. Frente a SQL, SPARQL permite navegar las relaciones del grafo de forma mucho más expresiva sin necesidad de múltiples JOINs. Por eso fue la elección natural para este proyecto.

\subsection*{Diapositiva: Docker}
Por último, Docker. Es la herramienta que nos permite empaquetar todos los servicios (backend Java, servicios Python, frontend) en contenedores ligeros. El usuario solo necesita instalar Docker y ejecutar un comando. La primera instalación descarga las imágenes, por lo que tarda unos minutos, pero las siguientes son casi inmediatas.

\section*{Sección 3: Comportamiento de la aplicación y manual de usuario}

\subsection*{Diapositiva: Comportamiento de la aplicación y manual de usuario}
Veamos ahora cómo se usa KnowGraph. \textit{(Señalar los puntos de la diapositiva)}

\subsection*{Diapositiva: Ejecución de consultas}
El flujo más básico es el de ejecutar una consulta SPARQL. El usuario escribe la consulta o la genera visualmente, pulsa ejecutar y obtiene una tabla paginada con los resultados, que puede descargar en CSV o JSON. Además, hay un chat con un asistente basado en inteligencia artificial para resolver dudas sobre SPARQL. \textit{(Señalar la captura de pantalla)}

\subsection*{Diapositiva: Generador visual de consultas SPARQL}
Para quienes no conocen SPARQL, disponemos de un “Query Wizard”. Se organiza en tres bloques: prefijos (para acortar URIs), datasets (opcional) y la construcción paso a paso de la consulta, donde se seleccionan variables, filtros, límites… Al finalizar, la consulta se copia automáticamente en el editor.

\subsection*{Diapositiva: Constructor de triplas RDF para inserciones}
De forma análoga, hay un asistente para construir inserciones de datos. Se definen triplas (sujeto, predicado, objeto) con ayuda de prefijos, se elige el grafo destino y se genera automáticamente la consulta INSERT. Pensado para usuarios que quieran poblar sus propios grafos sin escribir SPARQL manualmente.

\subsection*{Diapositiva: Cambio de proveedor de servicios SPARQL}
Una característica importante es la posibilidad de cambiar el motor que responde las consultas. Por defecto, usamos el servidor SPARQL local que incluye el proyecto, y eso nos da control total: cargar datos propios, exportarlos a cualquier formato, etc. Pero también podemos conectarnos a servicios públicos como Wikidata para lanzar consultas sobre sus grafos, aunque en ese caso se pierden las capacidades de carga y exportación local. Es un equilibrio entre autonomía y alcance.

\subsection*{Diapositiva: Visualización de datos en KnowGraph}
Tras obtener resultados, podemos visualizarlos de dos maneras.
\begin{itemize}
    \item La primera es como un \textbf{grafo interactivo}, en 2D o 3D, donde el usuario decide qué variables serán nodos y cuáles arcos.
    \item La segunda es mediante \textbf{reducción dimensional}: algoritmos como PCA, t‑SNE o UMAP proyectan los datos multidimensionales en dos o tres dimensiones, permitiendo detectar agrupaciones o patrones ocultos.
\end{itemize}
Todo esto se procesa bajo demanda y se cachea para no repetir cálculos costosos.

\subsection*{Diapositiva: Ejemplo de uso en la aplicación}
\textit{(Podría mostrarse una breve demo grabada o describir un flujo concreto).} Por ejemplo, cargar un CSV de películas, convertirlo a triplas RDF, ejecutar una consulta que relacione directores y géneros, y luego visualizar el resultado en un grafo 3D.

\section*{Sección 4: Documentación para desarrolladores}

\subsection*{Diapositiva: Documentación para desarrolladores}
Para los compañeros técnicos, entramos en la arquitectura del sistema.

\subsection*{Diapositiva: Introducción y arquitectura general}
KnowGraph sigue una arquitectura de microservicios. Tenemos un backend Java (que maneja el núcleo SPARQL con Apache Jena y Spring Boot), dos servicios Python (importer y visualizer) y un frontend React servido con Nginx. Todo está contenerizado con Docker y orquestado con Docker Compose. La comunicación se canaliza a través de Nginx, que actúa como proxy inverso.

\subsection*{Diapositiva: Backend Java}
El servicio Java utiliza Apache Jena para todas las operaciones sobre grafos RDF: carga, consulta, razonamiento… Spring Boot expone una API REST y Maven se encarga de construir el artefacto final.

\subsection*{Diapositiva: Backend Python}
Los dos servicios Python cumplen roles específicos.
\begin{itemize}
    \item El \textbf{Importer} es un Flask que recibe archivos (CSV, JSON, Excel, XML), los convierte en triplas RDF y las envía al backend Java.
    \item El \textbf{Visualizer}, también Flask, toma datos tabulares, aplica reducción de dimensionalidad con scikit‑learn y UMAP, y devuelve coordenadas listas para graficar.
\end{itemize}
Ambos utilizan Gunicorn como servidor de producción.

\subsection*{Diapositiva: Frontend}
El frontend está construido con React y TypeScript, usando Vite como bundler. Las rutas se gestionan con React Router y las peticiones a la API se redirigen mediante el proxy de Nginx. Proporciona la interfaz para todos los módulos: ejecución, construcción de consultas, carga de datos y visualización.

\subsection*{Diapositiva: Integración y flujos de trabajo}
El flujo típico es: el usuario interactúa con el frontend; este envía solicitudes HTTP a Nginx en el puerto 80; Nginx enruta según la ruta al servicio correspondiente (sparql‑backend, importer, visualizer). Docker Compose asegura que todos los contenedores estén en la misma red y se levanten en orden.

\subsection*{Diapositiva: Docker: contenerización y orquestación}
Cada servicio se define en el archivo \texttt{docker-compose.yaml}: imágenes, redes, volúmenes (como el volumen para los datos del Importer), healthchecks… Esto garantiza portabilidad total.

\subsection*{Diapositiva: Nginx como proxy inverso}
Nginx es la puerta de entrada. Unifica el acceso, sirve los estáticos y evita problemas de CORS. Además, hemos configurado timeouts amplios porque algunas operaciones de reducción dimensional pueden tardar varios segundos.

\subsection*{Diapositiva: Puntos fuertes, débiles y posibles mejoras}
\textit{(Aquí se disparan las tres diapositivas siguientes)}

\subsection*{Diapositiva: Puntos fuertes}
Entre los puntos fuertes destacan:
\begin{itemize}
    \item Despliegue sencillo con Docker.
    \item Potencia sobre datos locales (importación, transformación, consulta y exportación).
    \item Doble modalidad de visualización (grafo y reducción dimensional).
    \item Asistente basado en LLM que reduce la barrera de entrada a SPARQL.
\end{itemize}

\subsection*{Diapositiva: Puntos débiles}
Reconozcamos las limitaciones:
\begin{itemize}
    \item La dependencia de Docker puede ser una traba para usuarios no técnicos; la primera ejecución descarga muchas imágenes.
    \item La imagen del Visualizer no es la más ligera (usamos \texttt{python:slim} en vez de \texttt{alpine}).
    \item Al cambiar a un proveedor externo se pierden funcionalidades importantes.
    \item A nivel de código, hay deuda técnica y falta de patrones de diseño claros.
\end{itemize}

\subsection*{Diapositiva: Posibles mejoras}
En el futuro proponemos:
\begin{itemize}
    \item Una refactorización completa para mejorar la mantenibilidad.
    \item Ampliar el Importer a más formatos y fuentes remotas.
    \item Incorporar más algoritmos de visualización.
    \item Añadir una guía de ayuda interactiva.
    \item Dotar a la tabla de resultados de una barra de búsqueda y filtrado.
\end{itemize}
Todo esto haría la herramienta aún más robusta y accesible.


\section*{Bibliografía}
\subsection*{Diapositiva: Referencias bibliográficas}
Gran parte del trabajo se sustenta en el libro \textit{Learning SPARQL} de Bob DuCharme y en el artículo \textit{Knowledge Graphs} de Hogan et al. (2020). Además, se ha hecho un uso intensivo de la documentación oficial de Apache Jena, Spring Boot, Docker, Vite y Nginx, así como de los fundamentos de Linked Data propuestos por Tim Berners-Lee. Todo este material queda a su disposición para cualquier consulta.

\section*{Cierre}

\subsection*{Diapositiva: Cierre}
Y hasta aquí la presentación. Muchas gracias por su atención. Quedo a su disposición para cualquier pregunta.

\end{document}